\chapter*{Déclaration d'utilisation de l'intelligence artificielle}
\addcontentsline{toc}{chapter}{Déclaration d'utilisation de l'intelligence artificielle}

L'intelligence artificielle générative a été utilisée de manière ponctuelle au cours de
la réalisation de ce mémoire, principalement comme outil d'accompagnement dans certaines
étapes du travail.

Des outils tels que ChatGPT et Claude ont été sollicités pour faciliter certaines recherches
exploratoires, approfondir la compréhension de notions techniques, confronter des pistes de
réflexion ou encore améliorer la formulation et l'organisation de certains développements.
Ils ont également été utilisés pour la mise en forme de certains éléments techniques,
notamment en \LaTeX{}, ainsi que pour explorer des références ou des travaux susceptibles
d'enrichir la recherche bibliographique.

Ces outils ont par ailleurs constitué une aide à la programmation. Ils ont été mobilisés
pour l'écriture et la correction de portions de code relatives à l'apprentissage des modèles
de machine learning, à la construction des indicateurs de performance et aux représentations
graphiques, ainsi que pour la mise au point des scripts d'automatisation des lancements du
modèle ALM. Le code produit avec cette assistance a systématiquement été relu, testé et
validé, et la structure des traitements comme les choix d'implémentation relèvent du travail
mené dans le cadre de ce mémoire.

Ces outils n'ont toutefois pas constitué une source de données ou de résultats pour l'étude.
Les informations, références et éléments bibliographiques retenus ont été vérifiés à partir
de documents et de sources fiables. De la même manière, les choix de modélisation, les
développements effectués dans l'outil, les calculs, les résultats obtenus et leur analyse
relèvent du travail réalisé dans le cadre de ce mémoire.

\vspace{1.5cm}

\chapter*{Statement on the Use of Artificial Intelligence}
\addcontentsline{toc}{chapter}{Statement on the Use of Artificial Intelligence}

Generative artificial intelligence tools were used occasionally throughout the preparation
of this thesis, mainly as support during certain stages of the work.

Tools such as ChatGPT and Claude were used for exploratory research, to improve the
understanding of technical concepts, discuss methodological approaches, and refine the
wording and structure of certain sections. They were also used to assist with the formatting
of technical elements, particularly in \LaTeX{}, and to identify potentially relevant
references and publications.

These tools further provided assistance with programming. They supported the writing and
debugging of parts of the code relating to the training of the machine learning models, the
computation of performance metrics and the production of graphical outputs, as well as the
development of the scripts automating the ALM model runs. All code produced with such
assistance was reviewed, tested and validated, and both the structure of the processing
chain and the implementation choices are the result of the work carried out for this thesis.

These tools were not used as a source of data or actuarial results. The information,
references and bibliographical sources ultimately retained were checked against reliable
original sources. Similarly, the modelling choices, developments implemented in the tool,
calculations, numerical results and their interpretation are the result of the work carried
out for this thesis.